On September 7th 2025, there was a total lunar eclipse, visible from Australia, Asia, Africa and Europe. A detailed lunar eclipse figure, obtained from the website eclipsewise.com of eclipse predictions by Fred Espenak, is shown below. Based on the figure, respond to the following items.
\begin{description}
\item[(a)] Measure the inclination between the orbit of the Moon around the Earth and the Ecliptic plane. The orbit of the Moon is assumed to be circular.
\item[(b)] During the total eclipse phase, the moon appears red in colour. Which of the following factors are necessary to explain this red colour? Write their item letters in the answer sheet.
\begin{description}
\item[A)] Density of atmosphere
\item[B)] Albedo of Moon
\item[C)] Orbital radius of Moon
\item[D)] Albedo of Earth
\item[E)] Effective temperature of the Sun
\end{description}
\item[(c)] Measure the angular diameter of the Moon. Knowing the value of the radius of the Moon, calculate the Earth--Moon distance during the eclipse, in kilometres.
\item[(d)] Measure the angular diameter of Earth's umbra and the angular diameter of Earth's penumbra.
\item[(e)] What is the instantaneous angular velocity of the Moon relative to the observer, during the total eclipse? Give the answer in arcminutes per hour ($'/\mathrm h$).
\item[(f)] On 26 June 2029, a central total lunar eclipse will take place. This means that the centre of the disk of the Moon will pass through the axis of Earth's umbral shadow. If the angular velocity is assumed to be the same for the two eclipses, by what percentage will the 2029 total lunar eclipse last longer than the one in 2025 (percentage to be calculated of the totality duration in 2025)?
\item[(g)] Calculate the total duration of the eclipse in 2029 (from Penumbral first contact to Penumbral last contact).
\end{description}
\ioaafig{2025_TH_Q07-f1.pdf}
\ioaafig{2025_TH_Q07-f2.pdf}
